\scanpage{11}
\uline{Caractère} d'\uline{un groupe de} B.T. \uline{sur un corps local}.

Soit $V$ un anneau de val. discrète\inserted{\figure{\begin{tikzpicture}\node[draw,ellipse] {\uline{à corps résiduel parfait}};\end{tikzpicture}}}; pour ce qu'on va faire, on peut le supposer complet, \deleted{à corps résiduel algébriquement clos (\unsure{il me semble}\ \unsure{importe} semble l'hypothèse corps résiduel parfait; on peut alors, au besoin, \inserted{(quitte à} passer au complété de l'hensélisé strict)}. On suppose que le corps résiduel $k$ de $V$ est de car. $p\inserted{>}0$; le corps des fractions $K$ de $V$ est donc de car $p$ ou $0$. On \deleted{si} s'intéresse aux\inserted{$p$-}groupes de Barsotti-Tate sur $K$, notés $M=(M(\nu))_{\nu\geq0}$. On dit qu'un tel groupe \uline{satisfait à la condition de monodromie} s'il admet une suite de composition par des groupes de B.T., telle que les \inserted{$N=$}$\operatorname{Gr}^i(M)$ aient \uline{essentiellement bonne réduction}, i.e. qu'il existe pour chacun d'eux une extension finie \inserted{(séparable)} $K'$ de $K$ telle que \deleted{\unsure{sur $K$}} $N_{K'}$ ait bonne réduction sur le normalisé $S'$ de $S$ dans $K'$. Alors $N_{K'}$ provient \deleted{d'un} donc, d'après le th. de Tate (démontré pour l'instant seulement dans le cas d'\inserted{in}égales caractéristiques), d'un \deleted{\unsure{schéma en}} groupe de B.T. $N$ sur $S'$ déterminé \uline{à isomorphisme unique près}. Par suite, si $K'$ est prise galoisienne de groupe $G$, $G$ opère sur \deleted{\unsure{eux}} $N$ de façon compatible avec son action sur $S'$. - Bien entendu, on peut prendre $K'$ indépendant de $i$; on dira, si $K'$ marche pour $M$ satisfaisant la condition de monodromie, que $K'$ est un \uline{corps de bonne réduction virtuelle} pour $M$.

Soit $K_{\mathrm{BT}}(K;K')$ le groupe ``de Grothendieck'' \deleted{\unsure{déf}} défini\deleted{r} par les groupes de BT sur $K$ admettant $K'$ comme corps de bonne réduction virtuelle. Alors, si $K'/K$ est galoisienne de groupe de Galois $G$, on trouve un homomorphisme naturel
\diagram{\begin{tikzcd}[ampersand replacement=\&]
(1)\qquad K_{\mathrm{BT}}(K;K')\arrow[r]\&K_{\mathrm{BT}}(S';G),
\end{tikzcd}}
où le deuxième membre est le groupe de Grothendieck des groupes de B.T. sur $S'$ avec opérations de $G$ (compatibles aux opérations de $G$ sur $S'$). D'ailleurs, $S$ étant complet donc hensélien, $S'$ est local, de sorte que,
\mspage{10}
